% ---- pg-550 （书页 538） ----

理(此处每一对项目中的两个成员分别属于上述两个聚类).

同样, 平均连接算法的输入可以是距离或相似性, 而且这种方法也可以用于对象聚类或变量分类. 平均连接算法按式(12-12) 中的一般算法进行.
我们仍从距离阵 $\mathbf{D}=\{d_{ik}\}$ 中搜寻最接近(最相似)的对象开始, 假定找到的是 $U$ 和 $V$. 这两个对象合并, 组成聚类 $(UV)$. 对一般的聚集算法的步骤 3, 聚类 $(UV)$ 与另一聚类 $W$ 之间的距离由下式确定:

\phantomsection\label{eq:12-15}
\begin{equation}
d_{(UV)W}=\frac{\sum\limits_{i}\sum\limits_{k}d_{ik}}{N_{(UV)}N_{W}} \tag{12-15}
\end{equation}

其中 $d_{ik}$ 为聚类 $(UV)$ 中的第 $i$ 个对象与聚类 $W$ 中的第 $k$ 个对象之间的距离, $N_{(UV)}$ 和 $N_{W}$ 分别为 $(UV)$ 和 $W$ 中成员的个数.

\noindent\textbf{例 12.8}\quad (11 种语言的平均连接聚类)

将平均连接法用于例 12.4 中 11 种语言之间的距离, 所得出的树图示于\hyperref[fig:12.9]{图 12.9}.

\begin{mvfigure}
\begin{minipage}[t]{0.47\linewidth}\centering
\begin{tikzpicture}[x=0.60cm,y=1.9cm,line width=0.4pt]
% ---- 坐标轴 ----
\draw (-0.75,0) -- (-0.75,1.42);
\draw (-0.75,0.000) -- (-0.50,0.000);  \node[anchor=east,font=\fontsize{5}{6}\selectfont,inner sep=1pt] at (-0.80,0.000) {1.0};
\draw (-0.75,0.200) -- (-0.50,0.200);  \node[anchor=east,font=\fontsize{5}{6}\selectfont,inner sep=1pt] at (-0.80,0.200) {0.8};
\draw (-0.75,0.400) -- (-0.50,0.400);  \node[anchor=east,font=\fontsize{5}{6}\selectfont,inner sep=1pt] at (-0.80,0.400) {0.6};
\draw (-0.75,0.600) -- (-0.50,0.600);  \node[anchor=east,font=\fontsize{5}{6}\selectfont,inner sep=1pt] at (-0.80,0.600) {0.4};
\draw (-0.75,0.800) -- (-0.50,0.800);  \node[anchor=east,font=\fontsize{5}{6}\selectfont,inner sep=1pt] at (-0.80,0.800) {0.2};
\draw (-0.75,1.000) -- (-0.50,1.000);  \node[anchor=east,font=\fontsize{5}{6}\selectfont,inner sep=1pt] at (-0.80,1.000) {0};
\draw (-0.75,1.200) -- (-0.50,1.200);  \node[anchor=east,font=\fontsize{5}{6}\selectfont,inner sep=1pt] at (-0.80,1.200) {$-0.2$};
\draw (-0.75,1.400) -- (-0.50,1.400);  \node[anchor=east,font=\fontsize{5}{6}\selectfont,inner sep=1pt] at (-0.80,1.400) {$-0.4$};
\node[rotate=90,font=\fontsize{6}{7}\selectfont,inner sep=2pt] at (-1.45,0.70) {相似性(相关系数)};
% ---- 树图 ----
\draw (0.000,1.000) -- (0.000,0.350);
\draw (1.000,1.000) -- (1.000,0.350);
\draw (0.000,0.350) -- (1.000,0.350);
\draw (2.000,1.000) -- (2.000,0.950);
\draw (0.500,0.350) -- (0.500,0.950);
\draw (2.000,0.950) -- (0.500,0.950);
\draw (3.000,1.000) -- (3.000,0.500);
\draw (4.000,1.000) -- (4.000,0.500);
\draw (3.000,0.500) -- (4.000,0.500);
\draw (5.000,1.000) -- (5.000,0.570);
\draw (6.000,1.000) -- (6.000,0.570);
\draw (5.000,0.570) -- (6.000,0.570);
\draw (1.250,0.950) -- (1.250,1.050);
\draw (3.500,0.500) -- (3.500,1.050);
\draw (1.250,1.050) -- (3.500,1.050);
\draw (5.500,0.570) -- (5.500,0.990);
\draw (7.000,1.000) -- (7.000,0.990);
\draw (5.500,0.990) -- (7.000,0.990);
\draw (2.375,1.050) -- (2.375,1.210);
\draw (6.250,0.990) -- (6.250,1.210);
\draw (2.375,1.210) -- (6.250,1.210);
% ---- 叶子 ----
\foreach \i/\lb in {0/1,1/2,2/7,3/4,4/8,5/3,6/5,7/6} {
  \fill (\i,0) circle (1.0pt);
  \node[font=\fontsize{6}{7}\selectfont,inner sep=1pt,anchor=north,yshift=-2pt] at (\i,0) {\lb};
}
\node[font=\fontsize{6}{7}\selectfont,inner sep=4pt] at (3.5,-0.30) {变量};
\end{tikzpicture}
\phantomsection\label{fig:12.8}\mvcap{图 12.8\quad 八个公用事业公司变量之间\\相似性的完全连接树图}
\end{minipage}\hfill
\begin{minipage}[t]{0.47\linewidth}\centering
\begin{tikzpicture}[x=0.50cm,y=0.33cm,line width=0.4pt]
% ---- 坐标轴 ----
\draw (-0.75,0) -- (-0.75,10.3);
\draw (-0.75,0)  -- (-0.45,0);   \node[anchor=east,font=\fontsize{5}{6}\selectfont,inner sep=1pt] at (-0.80,0)  {0};
\draw (-0.75,2)  -- (-0.45,2);   \node[anchor=east,font=\fontsize{5}{6}\selectfont,inner sep=1pt] at (-0.80,2)  {2};
\draw (-0.75,4)  -- (-0.45,4);   \node[anchor=east,font=\fontsize{5}{6}\selectfont,inner sep=1pt] at (-0.80,4)  {4};
\draw (-0.75,6)  -- (-0.45,6);   \node[anchor=east,font=\fontsize{5}{6}\selectfont,inner sep=1pt] at (-0.80,6)  {6};
\draw (-0.75,8)  -- (-0.45,8);   \node[anchor=east,font=\fontsize{5}{6}\selectfont,inner sep=1pt] at (-0.80,8)  {8};
\draw (-0.75,10) -- (-0.45,10);  \node[anchor=east,font=\fontsize{5}{6}\selectfont,inner sep=1pt] at (-0.80,10) {10};
\draw (-0.75,1) -- (-0.55,1); \draw (-0.75,3) -- (-0.55,3);
\draw (-0.75,5) -- (-0.55,5); \draw (-0.75,7) -- (-0.55,7);
\draw (-0.75,9) -- (-0.55,9);
\node[rotate=90,font=\fontsize{6}{7}\selectfont,inner sep=2pt] at (-1.75,5) {距离};
% ---- 树图 ----
\draw (1.000,0.000) -- (1.000,1.000);
\draw (2.000,0.000) -- (2.000,1.000);
\draw (1.000,1.000) -- (2.000,1.000);
\draw (0.000,0.000) -- (0.000,2.000);
\draw (1.500,1.000) -- (1.500,2.000);
\draw (0.000,2.000) -- (1.500,2.000);
\draw (5.000,0.000) -- (5.000,1.000);
\draw (6.000,0.000) -- (6.000,1.000);
\draw (5.000,1.000) -- (6.000,1.000);
\draw (5.500,1.000) -- (5.500,1.500);
\draw (7.000,0.000) -- (7.000,1.500);
\draw (5.500,1.500) -- (7.000,1.500);
\draw (6.250,1.500) -- (6.250,4.000);
\draw (8.000,0.000) -- (8.000,4.000);
\draw (6.250,4.000) -- (8.000,4.000);
\draw (0.750,2.000) -- (0.750,5.000);
\draw (3.000,0.000) -- (3.000,5.000);
\draw (0.750,5.000) -- (3.000,5.000);
\draw (1.875,5.000) -- (1.875,6.000);
\draw (4.000,0.000) -- (4.000,6.000);
\draw (1.875,6.000) -- (4.000,6.000);
\draw (2.938,6.000) -- (2.938,7.000);
\draw (7.125,4.000) -- (7.125,7.000);
\draw (2.938,7.000) -- (7.125,7.000);
\draw (9.000,0.000) -- (9.000,8.000);
\draw (10.000,0.000) -- (10.000,8.000);
\draw (9.000,8.000) -- (10.000,8.000);
\draw (5.031,7.000) -- (5.031,9.000);
\draw (9.500,8.000) -- (9.500,9.000);
\draw (5.031,9.000) -- (9.500,9.000);
% ---- 叶子 ----
\foreach \i/\lb in {0/E,1/N,2/Da,3/G,4/Du,5/Fr,6/I,7/Sp,8/P,9/H,10/Fi} {
  \fill (\i,0) circle (1.0pt);
  \node[font=\fontsize{6}{7}\selectfont,inner sep=1pt,anchor=north,yshift=-2pt] at (\i,0) {\lb};
}
\node[font=\fontsize{6}{7}\selectfont,inner sep=4pt] at (5,-2.5) {语言};
\end{tikzpicture}
\phantomsection\label{fig:12.9}\mvcap{图 12.9\quad 11 种语言中前 10 个数词之间\\距离的平均连接树图}
\end{minipage}
\end{mvfigure}

将\hyperref[fig:12.9]{图 12.9} 中的树图与单连接树图(\hyperref[fig:12.4]{图 12.4}), 以及完全连接树图(\hyperref[fig:12.7]{图 12.7})相比较, 就会发现平均连接产生的结构与完全连接的结构非常一致.
然而由于每一种情况下距离的定义不同, 合并在不同的水平产生是不奇怪的. \hfill$\blacksquare$

\noindent\textbf{例 12.9}\quad (公用事业公司的平均连接聚类)

将平均连接算法应用于 22 个公用事业公司的欧氏距离(见\hyperref[tab:12.6]{表 12.6}), 便产生\hyperref[fig:12.10]{图 12.10} 中的树图.

\begin{mvfigure}
\begin{tikzpicture}[x=0.45cm,y=1.05cm,line width=0.4pt]
% ---- 坐标轴 ----
\draw (-0.90,0) -- (-0.90,4.25);
\draw (-0.90,0) -- (-0.62,0);  \node[anchor=east,font=\fontsize{5}{6}\selectfont,inner sep=1pt] at (-0.95,0) {0};
\draw (-0.90,1) -- (-0.62,1);  \node[anchor=east,font=\fontsize{5}{6}\selectfont,inner sep=1pt] at (-0.95,1) {1};
\draw (-0.90,2) -- (-0.62,2);  \node[anchor=east,font=\fontsize{5}{6}\selectfont,inner sep=1pt] at (-0.95,2) {2};
\draw (-0.90,3) -- (-0.62,3);  \node[anchor=east,font=\fontsize{5}{6}\selectfont,inner sep=1pt] at (-0.95,3) {3};
\draw (-0.90,4) -- (-0.62,4);  \node[anchor=east,font=\fontsize{5}{6}\selectfont,inner sep=1pt] at (-0.95,4) {4};
\draw (-0.90,0.5) -- (-0.75,0.5); \draw (-0.90,1.5) -- (-0.75,1.5);
\draw (-0.90,2.5) -- (-0.75,2.5); \draw (-0.90,3.5) -- (-0.75,3.5);
\node[rotate=90,font=\fontsize{6}{7}\selectfont,inner sep=2pt] at (-1.70,2.1) {距离};
% ---- 树图（横线） ----
\draw (0.000,1.840) -- (0.970,1.840);
\draw (0.490,1.960) -- (2.480,1.960);
\draw (1.980,1.840) -- (2.980,1.840);
\draw (8.080,1.360) -- (9.090,1.360);
\draw (8.570,1.840) -- (10.510,1.840);
\draw (10.070,1.720) -- (11.130,1.720);
\draw (7.030,2.280) -- (9.450,2.280);
\draw (13.130,1.280) -- (14.110,1.280);
\draw (13.640,2.000) -- (15.090,2.000);
\draw (14.420,2.160) -- (16.120,2.160);
\draw (12.100,2.460) -- (15.270,2.460);
\draw (18.090,2.220) -- (19.110,2.220);
\draw (17.090,3.000) -- (18.740,3.000);
\draw (1.490,2.770) -- (6.020,2.770);
\draw (2.930,2.850) -- (8.250,2.850);
\draw (5.550,2.920) -- (13.720,2.920);
\draw (9.650,3.180) -- (17.990,3.180);
\draw (13.750,3.620) -- (20.000,3.620);
\draw (16.950,4.120) -- (20.980,4.120);
\draw (4.020,2.680) -- (4.990,2.680);
% ---- 树图（竖线） ----
\draw (0.460,1.840) -- (0.460,1.960);
\draw (2.480,1.840) -- (2.480,1.960);
\draw (4.490,2.680) -- (4.490,2.770);
\draw (8.570,1.360) -- (8.570,1.840);
\draw (10.600,1.720) -- (10.600,1.840);
\draw (9.450,1.840) -- (9.450,2.280);
\draw (8.250,2.280) -- (8.250,2.850);
\draw (13.640,1.280) -- (13.640,2.000);
\draw (14.420,2.000) -- (14.420,2.160);
\draw (15.270,2.160) -- (15.270,2.460);
\draw (18.580,2.220) -- (18.580,3.000);
\draw (17.990,3.000) -- (17.990,3.180);
\draw (2.930,2.770) -- (2.930,2.850);
\draw (5.550,2.850) -- (5.550,2.920);
\draw (9.650,2.920) -- (9.650,3.180);
\draw (13.750,3.180) -- (13.750,3.620);
\draw (16.950,3.620) -- (16.950,4.120);
\draw (1.490,1.960) -- (1.490,2.770);
\draw (0.000,0) -- (0.000,1.840);
\draw (1.000,0) -- (1.000,1.840);
\draw (2.000,0) -- (2.000,1.840);
\draw (3.000,0) -- (3.000,1.840);
\draw (4.000,0) -- (4.000,2.680);
\draw (5.000,0) -- (5.000,2.680);
\draw (6.000,0) -- (6.000,2.770);
\draw (7.000,0) -- (7.000,2.280);
\draw (8.000,0) -- (8.000,1.360);
\draw (9.000,0) -- (9.000,1.360);
\draw (10.000,0) -- (10.000,1.720);
\draw (11.000,0) -- (11.000,1.720);
\draw (12.000,0) -- (12.000,2.460);
\draw (13.000,0) -- (13.000,1.280);
\draw (14.000,0) -- (14.000,1.280);
\draw (15.000,0) -- (15.000,2.000);
\draw (16.000,0) -- (16.000,2.160);
\draw (17.000,0) -- (17.000,3.000);
\draw (18.000,0) -- (18.000,2.220);
\draw (19.000,0) -- (19.000,2.220);
\draw (20.000,0) -- (20.000,3.620);
\draw (21.000,0) -- (21.000,4.120);
% ---- 叶子 ----
\foreach \i/\lb in {0/1,1/18,2/19,3/14,4/9,5/3,6/6,7/22,8/10,9/13,10/20,11/4,12/7,13/12,14/21,15/15,16/2,17/11,18/16,19/8,20/5,21/17} {
  \fill (\i,0) circle (0.8pt);
  \node[font=\fontsize{5}{6}\selectfont,inner sep=0.6pt,anchor=north,yshift=-2pt] at (\i,0) {\lb};
}
\node[font=\fontsize{6}{7}\selectfont,inner sep=6pt] at (10.5,-1.25) {公用事业公司};
\end{tikzpicture}
\phantomsection\label{fig:12.10}\mvcap{图 12.10\quad 22 个公用事业公司间距离的平均连接树图}
\end{mvfigure}
